\scanpage{48}
\mspage{\handwritten{1}}

% printed form
\deleted{J0369CNU}

\handwritten{15.1.}

\handwritten{$X$ \deleted{non singulière}\inserted{lisse type $g,\unsure{\nu|v}$} (sur $\overline{\mathbb Q}$)}

\handwritten{$Y$ stable non lisse type $(g,\nu)$ avec propagateur normal $N$, d'où un groupe fondamental \deleted{modifié}\inserted{ext.} $\pi_1\deleted{\unsure{\langle|(}}[Y,N\deleted{\unsure{\rangle|)}}]$.}

\handwritten{$X\underset{T}{\overset{\unsure{\varphi|f|\tau}}{\simeq}}Y$ un isomorphisme de Teichmüller.}

\handwritten{\uline{NB} Cet isomorphisme induit un isom de groupes extérieurs \deleted{\unsure{(et est dé|Ceci est de}} (et est déterminé par lui)
\[
\pi_1(X)\longrightarrow\pi_1(\deleted{\unsure{\langle|(}}[Y,N\deleted{\unsure{\rangle|)}}]).
\]}

\handwritten{\deleted{\unsure{On|O}}Quand on se donne un ``pt base sur $X$'' (ordinaire, ou à l'infini), et un sur $[Y,N]$ (idem -- et de plus distinct des pts singuliers de $Y$) on a une notion d'isom de Teichmüller \unsure{respectant|respect. nt} les pts marqués, \deleted{\unsure{et dét. un isom|et est un isom}}\inserted{(on suppose} qu'ils sont tous deux à distance finie, ou tous deux à l'infini) et on trouve alors un isom de groupes
\[
\text{(pas ext. !)}\qquad \pi_1(X,\unsure{x|\bullet})\longrightarrow\pi_1([Y,N],y)
\]}

\handwritten{En tout état de cause, \deleted{\unsure{l'iso|il}} l'isom extérieur déjà n'est pas compatible avec les actions ext. de $\mathrm{Gal}\,(\overline{\mathbb Q}/K)$, $K$ une \inserted{ss} extension finie de $\overline{\mathbb Q}/\mathbb Q$ assez grande.}

\handwritten{\uline{NB} Mettre en évidence \uline{trois} classes de conjugaison d'\deleted{éléments} hom. $\widehat{\mathbb Z}(1)\longrightarrow\pi_1(X)$ (deux seulement si $\nu\unsure{\geq|=|\ne|>}0$, l'une des trois étant les ``lacets'').}

% archive sheet number: 50
